\documentclass[../main.tex]{subfiles}
\graphicspath{{\subfix{../img/}}{\subfix{../architectures images/}}}
\begin{document}
\chapter{Approche de D\'{e}veloppement et Technologies Utilis\'{e}es}

\section{Introduction}
Ce chapitre explique comment le projet TrustGym a \'{e}t\'{e} organis\'{e} et construit : m\'{e}thodologie Scrum, pile technologique, architectures logique et physique, et environnement de travail.

%──────────────────────────────────────────────
\section{M\'{e}thodologie}
%──────────────────────────────────────────────

\subsection{M\'{e}thodologie Agile et Scrum}

Face \`{a} des besoins \'{e}volutifs et des retours partenaires continus, nous avons adopt\'{e} \textbf{Scrum} \cite{schwaber2020scrum, scrum, dev_methodology} : sprints de 2 \`{a} 4 semaines, daily standups (15~min), sprint planning et sprint review. M.~Bilel Taktak (CTO) a tenu le r\^{o}le de Product Owner et M.~Talel Zouari le suivi p\'{e}dagogique. Le projet s'est d\'{e}roul\'{e} en \textbf{sept sprints}.

\begin{figure}[H]
  \centering
  \includegraphics[width=0.85\textwidth]{scrum logo.png}
  \caption{Processus Scrum appliqu\'{e} au projet TrustGym}
  \label{fig:agile_scrum}
\end{figure}

%──────────────────────────────────────────────
\section{Technologies Utilis\'{e}es}
%──────────────────────────────────────────────

Le tableau ci-dessous r\'{e}capitule la pile technologique de TrustGym.

\begin{figure}[H]
  \centering
  \includegraphics[width=0.85\textwidth]{Technologie utilisee.png}
  \caption{Pile technologique de la plateforme TrustGym}
  \label{fig:tech_stack}
\end{figure}

\begin{longtable}{|p{3.5cm}|p{3cm}|p{7cm}|}
\hline
\textbf{Cat\'{e}gorie} & \textbf{Outil / Version} & \textbf{R\^{o}le dans le projet} \\
\hline
\endfirsthead
\hline
\endfoot
Frontend Web & Next.js 15, TanStack Query 5, Tailwind CSS & Dashboard admin/caissier avec \acrshort{ssr}, App Router, i18n (fr/ar/en) \cite{nextjs, tanstack, tailwind}. \\
\hline
Frontend Mobile & React Native, Expo \acrshort{sdk} 54, NativeWind & Application adh\'{e}rents iOS/Android, Expo Router \cite{reactnative, expo, nativewind}. \\
\hline
Backend & NestJS 10, TypeORM, Passport.js & API \acrshort{rest} modulaire (50+ modules), double \acrshort{jwt} (admin vs membre) \cite{nestjs, typeorm, passport}. \\
\hline
Base de donn\'{e}es & PostgreSQL 16 + pgvector & Stockage relationnel + recherche vectorielle pour le \acrshort{rag} \cite{postgre, pgvector}. \\
\hline
Temps r\'{e}el & Socket.IO, PG LISTEN/NOTIFY & Diffusion des \'{e}v\'{e}nements d'acc\`{e}s et d'occupation sur le namespace \texttt{/attendance} \cite{socketio}. \\
\hline
Intelligence Artificielle & OpenAI \texttt{gpt-4o-mini} + \texttt{text-embedding-3-small} & Pipeline \acrshort{rag} : embedding $\rightarrow$ pgvector $\rightarrow$ reranking $\rightarrow$ g\'{e}n\'{e}ration \cite{openai}. \\
\hline
Mat\'{e}riel & ZKTeco \acrshort{sdk} (TCP 4370), Hikvision ISAPI (\acrshort{https}) & Contr\^{o}le d'acc\`{e}s physique via Machine Handler (PM2, ngrok) \cite{zkteco, hikvision, pm2}. \\
\hline
Infra / Autres & VPS OVH, Nginx, Vercel, AWS S3, Nodemailer, Helmet.js & H\'{e}bergement, reverse proxy, stockage fichiers, emails transactionnels, s\'{e}curit\'{e} HTTP \cite{aws, helmet}. \\
\hline
\caption{Pile technologique de TrustGym}
\end{longtable}

%──────────────────────────────────────────────
\section{Architecture Syst\`{e}me}
%──────────────────────────────────────────────

L'architecture TrustGym s\'{e}pare clairement les responsabilit\'{e}s en quatre sous-projets : un par besoin m\'{e}tier, aucun ne fait le travail d'un autre.

\subsection{Architecture Logique}

La figure~\ref{fig:arch_logique} montre comment les quatre sous-projets communiquent entre eux.

\begin{itemize}
    \item \textbf{trustgym-api} : L'\acrshort{api} \acrshort{rest} centrale d\'{e}velopp\'{e}e avec NestJS constitue le c\oe{}ur du syst\`{e}me. Elle expose l'ensemble des endpoints consomm\'{e}s par les clients et g\`{e}re la logique m\'{e}tier, l'authentification, et la persistance des donn\'{e}es via TypeORM et PostgreSQL 16.

    \item \textbf{trustgym-app} : Le tableau de bord d'administration d\'{e}velopp\'{e} avec Next.js 15, destin\'{e} aux administrateurs et aux caissiers. L'authentification est g\'{e}r\'{e}e via NextAuth.js \cite{nextauth}, l'internationalisation via next-intl \cite{nextintl}, et les formulaires via React Hook Form \cite{rhf} avec validation Yup \cite{yup}. Il consomme l'\acrshort{api} \acrshort{rest} et se connecte au namespace WebSocket \texttt{/attendance} pour recevoir les mises \`{a} jour d'occupation en temps r\'{e}el.

    \item \textbf{trustgym-mobile} : L'application mobile d\'{e}velopp\'{e}e avec React Native et Expo \acrshort{sdk} 54, destin\'{e}e aux membres de la salle. Elle consomme \'{e}galement l'\acrshort{api} \acrshort{rest} avec un syst\`{e}me d'authentification \acrshort{jwt} s\'{e}par\'{e} via \texttt{MemberJwtGuard}.

    \item \textbf{trustgym-machine-handler} : Le pont mat\'{e}riel compos\'{e} d'un service Python (port 8000) communiquant directement avec les appareils ZKTeco/Hikvision et d'un service NestJS (port 8080) transmettant les \'{e}v\'{e}nements \`{a} l'\acrshort{api} centrale.
\end{itemize}

\begin{figure}[H]
  \centering
  \includegraphics[width=\textwidth]{ARCHITECTURE LOGIQUE.png}
  \caption{Architecture logique de la plateforme TrustGym}
  \label{fig:arch_logique}
\end{figure}

Le temps r\'{e}el repose sur Socket.IO (namespace \texttt{/attendance}) et le m\'{e}canisme natif PostgreSQL LISTEN/NOTIFY : le service \texttt{PgListenerService} \'{e}coute le canal \texttt{access\_log\_channel} et propage instantan\'{e}ment les \'{e}v\'{e}nements d'entr\'{e}e/sortie \`{a} tous les clients connect\'{e}s.

\subsection{Architecture Physique}

Ct\'{e} infrastructure : un \acrshort{vps} OVH sous Ubuntu, avec Nginx en reverse proxy devant les services backend. La figure~\ref{fig:arch_physique} d\'{e}taille la configuration de d\'{e}ploiement.

\begin{figure}[H]
  \centering
  \includegraphics[width=0.9\textwidth]{ARCHITECTURE PHYSIQUE.png}
  \caption{Architecture physique de d\'{e}ploiement de TrustGym}
  \label{fig:arch_physique}
\end{figure}

Le frontend est sur Vercel (\acrshort{cdn} Edge). L'application mobile se distribue via EAS Build. L'\acrshort{api} backend consomme AWS S3, Nodemailer et OpenAI comme services externes.

\subsection{Architecture Modulaire}

Chaque domaine fonctionnel de l'\acrshort{api} NestJS vit dans son propre module, selon le patron :

\begin{center}
    \textbf{Module} $\rightarrow$ \textbf{Controller} $\rightarrow$ \textbf{Service} $\rightarrow$ \textbf{Entity (TypeORM)}
\end{center}

Plus de 50 modules de domaine s'enregistrent dans \texttt{AppModule}. La s\'{e}curit\'{e} est assur\'{e}e par une cha\^{i}ne de guards :

\begin{itemize}
    \item \textbf{JwtAuthGuard} : V\'{e}rification du token \acrshort{jwt} pour les administrateurs, caissiers et super-administrateurs.
    \item \textbf{MemberJwtGuard} : V\'{e}rification du token \acrshort{jwt} sp\'{e}cifique aux membres de l'application mobile.
    \item \textbf{RolesGuard} : Contr\^{o}le d'acc\`{e}s bas\'{e} sur les r\^{o}les (\acrshort{rbac}) via le d\'{e}corateur \texttt{@Roles()}.
    \item \textbf{ModuleGuard} : Contr\^{o}le d'acc\`{e}s aux modules \acrshort{saas} via le d\'{e}corateur \texttt{@RequiresModule()}, permettant d'activer ou d\'{e}sactiver des fonctionnalit\'{e}s par partenaire (feature flags).
\end{itemize}

L'alias \texttt{@/*} pointe vers \texttt{src/*} dans tous les sous-projets --- un d\'{e}tail simple qui \'{e}vite les imports relatifs inextricables.

\subsection{Architecture Multi-Tenant}

TrustGym sert plusieurs clients sur la m\^{e}me instance, avec une isolation stricte des donn\'{e}es. La hi\'{e}rarchie est \`{a} trois niveaux :

\begin{itemize}
    \item \textbf{Plateforme} (TrustGym \acrshort{saas}) : instance unique h\'{e}bergeant l'ensemble des locataires.
    \item \textbf{Partenaire} (\texttt{Partner}) : le locataire de premier niveau, repr\'{e}sentant une franchise ou un propri\'{e}taire de cha\^{i}ne de salles. Il d\'{e}tient la configuration fiscale (\texttt{taxIdentificationNumber}, \texttt{tva}, \texttt{timbre}), les param\`{e}tres de cotisation annuelle et les formules d'abonnement partag\'{e}es.
    \item \textbf{Salle} (\texttt{Gym}) : le locataire de second niveau, repr\'{e}sentant un \'{e}tablissement physique. Chaque salle poss\`{e}de ses propres membres, coachs, utilisateurs, portes, stocks et \'{e}v\'{e}nements d'acc\`{e}s.
\end{itemize}

Le champ \texttt{partnerId} pr\'{e}sent sur la majorit\'{e} des entit\'{e}s est le verrou de cette isolation : chaque requ\^{e}te est filtr\'{e}e automatiquement, ce qui rend impossible l'acc\`{e}s aux donn\'{e}es d'un autre partenaire, m\^{e}me accidentellement.

\subsection{Architecture de S\'{e}curit\'{e}}

La s\'{e}curit\'{e} de TrustGym se d\'{e}compose en quatre m\'{e}canismes distincts :

\subsubsection{Double Syst\`{e}me JWT}

Deux cha\^{i}nes \acrshort{jwt} \cite{jwt_rfc7519} coexistent, avec le m\^{e}me secret mais des payloads diff\'{e}rents. Le champ \texttt{type} est le discriminant :

\begin{itemize}
    \item \textbf{Canal Admin/Caissier} : le token ne contient pas de champ \texttt{type}. Payload : \texttt{\{id, role, partnerId, gymId, sdk\}}. Valid\'{e} par \texttt{JwtStrategy} et prot\'{e}g\'{e} par \texttt{JwtAuthGuard}.
    \item \textbf{Canal Membre (mobile)} : le token contient \texttt{type: 'member'}. Payload : \texttt{\{id, type: 'member'\}}. Valid\'{e} par \texttt{JwtMemberStrategy} et prot\'{e}g\'{e} par \texttt{MemberJwtGuard}.
\end{itemize}

Chaque guard rejette les tokens de l'autre canal : un membre ne peut pas utiliser un token admin, et vice versa.

\subsubsection{Cha\^{i}ne de Gardes}

Les routes prot\'{e}g\'{e}es s'expriment via des d\'{e}corateurs NestJS, par exemple :

\begin{center}
    \texttt{@UseGuards(JwtAuthGuard, RolesGuard)} $\rightarrow$ \texttt{@Roles(RoleEnum.Admin)}
\end{center}

\begin{itemize}
    \item \textbf{JwtAuthGuard} : valide le token et injecte l'utilisateur dans le contexte de la requ\^{e}te via \texttt{AsyncLocalStorage} pour la tra\c{c}abilit\'{e}.
    \item \textbf{RolesGuard} : v\'{e}rifie que le r\^{o}le de l'utilisateur correspond aux r\^{o}les autoris\'{e}s par le d\'{e}corateur \texttt{@Roles()}.
    \item \textbf{ModuleGuard} : v\'{e}rifie que le module \acrshort{saas} demand\'{e} est activ\'{e} pour le partenaire via \texttt{SaasModuleEnum}.
\end{itemize}

\subsubsection{Authentification Machine-\`{a}-Machine}

La communication Machine Handler $\leftrightarrow$ \acrshort{api} utilise des cl\'{e}s d\'{e}di\'{e}es, diff\'{e}rentes selon la direction :

\begin{itemize}
    \item \textbf{Cloud $\rightarrow$ Bridge} : en-t\^{e}te \texttt{x-bridge-key} avec comparaison temporellement s\^{u}re (\textit{timing-safe}) pour \'{e}viter les attaques par canal auxiliaire.
    \item \textbf{Bridge $\rightarrow$ Python} : en-t\^{e}te \texttt{x-api-key} avec v\'{e}rification HMAC (\texttt{compare\_digest}).
\end{itemize}

\subsubsection{Impersonation SuperAdmin}

Le SuperAdmin peut se connecter en tant qu'un autre utilisateur pour le d\'{e}bogage et le support. Le token \acrshort{jwt} porte les champs \texttt{isImpersonation}, \texttt{impersonatedBy} et \texttt{impersonatorEmail}. Quelle que soit l'action r\'{e}alis\'{e}e, la piste d'audit conserve l'identit\'{e} r\'{e}elle de l'op\'{e}rateur --- personne ne peut se cacher derrière l'impersonation.

\subsection{Internationalisation (i18n)}

Le tableau de bord supporte trois langues via \texttt{next-intl} \cite{nextintl} :

\begin{itemize}
    \item \textbf{Fran\c{c}ais} (fr) : langue par d\'{e}faut.
    \item \textbf{Arabe} (ar) : d\'{e}clenche automatiquement la mise en page \acrshort{rtl} (\textit{Right-to-Left}).
    \item \textbf{Anglais} (en) : support complet.
\end{itemize}

Le middleware Next.js d\'{e}tecte la langue du navigateur et redirige vers la bonne locale (\texttt{/[locale]/page}). L'utilisateur peut aussi basculer manuellement via le \texttt{LanguageToggle} dans la barre de navigation.

%──────────────────────────────────────────────
\section{Environnement de Travail et Outils}
%──────────────────────────────────────────────

\subsection{Outils de D\'{e}veloppement}

\begin{itemize}
    \item \textbf{IDE} : Visual Studio Code \cite{vscode} avec extensions TypeScript/NestJS.
    \item \textbf{Mod\'{e}lisation UML} : StarUML (diagrammes UC, classes, s\'{e}quence).
    \item \textbf{Tests \acrshort{api}} : Postman \cite{postman} + Swagger auto-g\'{e}n\'{e}r\'{e} \`{a} \texttt{/api/doc} \cite{swagger}.
    \item \textbf{Gestion de projet} : Linear (sprint backlog, cycles automatis\'{e}s) ; Google Meet (daily standups et reviews).
    \item \textbf{Contr\^{o}le de version} : Git + GitHub (pull requests, revues de code) \cite{git}.
\end{itemize}

\subsection{Mat\'{e}riel de D\'{e}veloppement}

Le d\'{e}veloppement a \'{e}t\'{e} r\'{e}alis\'{e} sur la configuration suivante :
\begin{itemize}
    \item \textbf{Syst\`{e}me d'exploitation} : Windows 11 Pro
    \item \textbf{Serveur de base de donn\'{e}es} : PostgreSQL 16 (port 5433)
    \item \textbf{Mat\'{e}riel de test} : Terminaux ZKTeco C3-200 et cam\'{e}ras Hikvision SC705 pour le contr\^{o}le d'acc\`{e}s
    \item \textbf{Appareils mobiles} : Simulateurs iOS/Android via Expo Go, tests sur appareils physiques
\end{itemize}

\subsection{Workflow de D\'{e}veloppement et CI/CD}

Le workflow de d\'{e}veloppement est lin\'{e}aire, du sprint backlog jusqu'au d\'{e}ploiement, avec GitHub Actions comme colonne vert\'{e}brale du \acrshort{cicd} :

\begin{enumerate}
    \item \textbf{Planification} : D\'{e}finition des user stories dans Linear lors du Sprint Planning, avec attribution des priorit\'{e}s et estimation des charges.
    \item \textbf{D\'{e}veloppement} : Cr\'{e}ation de branches feature sur Git, d\'{e}veloppement local avec hot reload (NestJS, Next.js, Expo).
    \item \textbf{Tests} : Validation des endpoints via Swagger/Postman, tests manuels de l'interface, tests unitaires Jest.
    \item \textbf{Revue de code} : Pull requests sur GitHub avec revue par l'encadrant entreprise.
    \item \textbf{Int\'{e}gration continue} : GitHub Actions ex\'{e}cute automatiquement les \'{e}tapes de lint, test et build \`{a} chaque push.
    \item \textbf{D\'{e}ploiement} : Le frontend est d\'{e}ploy\'{e} sur Vercel (Path A), le backend sur le \acrshort{vps} OVH via SSH et PM2 (Path B). Des notifications Discord informent l'\'{e}quipe du statut de chaque d\'{e}ploiement.
\end{enumerate}

\begin{figure}[H]
  \centering
  \includegraphics[width=\textwidth]{ARCHITECTURE CICD.png}
  \caption{Pipeline CI/CD via GitHub Actions}
  \label{fig:cicd_pipeline}
\end{figure}

\section{Conclusion}
Ce chapitre a pos\'{e} les bases de la construction de TrustGym : pourquoi Scrum, quelles technologies et pourquoi, comment l'architecture est decoup\'{e}e. Le chapitre suivant passe \`{a} ce que le syst\`{e}me doit faire --- la sp\'{e}cification des exigences.

\end{document}
